\documentclass[../main.tex]{subfiles}
\graphicspath{{\subfix{../images/}}}
\begin{document}

It is generally known, that too much or not enough body weight can have consequences on our health.
A weight below the ideal weight indicates an under-supply of vital nutrients, which expresses itself in deficiency syndromes,
like tiredness, in a decrease of power, susceptibility to infections and so on.
Overweight indicates an oversupply with energy rich nutrients and/or a lack of movement.
Severe diseases of the heart--circulation system, metabolism problems like diabetes type 2, up to the metabolic syndrome are the consequences.

But what's a healthy body weight? At which point is it considered overweight or underweight?
From a medical viewpoint, there's no consensus, how much the ``normal body weight'' of a human being should be.
For  a long time, the research of a French surgeon Broca (1824-1880) was the basis to calculate the normal, ideal and overweight:

\noindent
\begin{tabular}{lll}
  Normal weight & = Height in cm &-- 100, ex: 175 cm -- 100 = 75 kg \\
  Ideal weight for men & = Normal weight &-- 10\% \\
  Ideal weight for women & = Normal weight &-- 5\% \\
  Overweight & = Normal weight &+ 10\% \\
\end{tabular}

The big disadvantage of the Broca formula is the fact, that it is only applicable to people with a height between 1.60 m and 1.80 m (5ft3 -- 5ft11).

\section{Body Mass Index or BMI}

Nowadays, the body mass index (BMI)\index{bmi}\index{body mass index BMI} is the standard to judge the body weight.
The necessary metrics of body weight, and height are easy to evaluate exactly.

The formula to calculate the BMI with an example:

\vspace{2mm}
\begin{center}
  \begin{fminipage}{10cm}
  $BMI = \frac{Body weight [kg]}{Height [m]^2} = 703 * \frac{Body weight [lbs]}{Height [in]}
  = \frac{75 kg}{(1.75 m)^2} = 24.5$
  \end{fminipage}
\end{center}

With the help of the calculated BMI value and the following table, it can be determined if there's under--, normal-- or overweight and the corresponding health risk.

\begin{table}[htb!]
  \centering
  \begin{tabular}{lll}
    \textbf{BMI [kg/m$^2$]} & \textbf{Classification} & \textbf{Health risk} \\
    \hline
    $<$ 18.5 & Underweight & low \\
    18.5 -- 24.9 & Normal weight & normal \\
    25.0 --29.9 & Overweight & increased \\
    30.0 -- 34.9 & Adiposity class 1 & strongly increased \\
    35.0 -- 39.9 & Adiposity class 2 & very strongly increased \\
    $\geq$ 40 & Adiposity class 3 & extremely  increased \\
  \end{tabular}
\end{table}

\section[Waist--hip--ratio (WHR)]{Waist--hip--ratio (WHR) or How is the Fat Distributed}
In order to accurately gauge the health risks in connection with overweight, the distribution of the excess fat pads also plays a role.
This can be determined by measuring the circumference of the body.
From the quotient of waist circumference to hip circumference, the so called waist--hip--ratio can be determined.


In order to be measured, the patient has to be standing up:
\begin{itemize}
\item the circumference of the waist, between the upper edge of the hip and the first rip
  \item the circumference of the hip at the largest place
  \end{itemize}

  Formula to calculate the WHR, with the example for a woman:

  \vspace{2mm}

\begin{center}
\begin{fminipage}{8cm}
$WHR = \frac{circumference\ waist [cm/in]\footnote{As the WHR is a ratio of measures, it doesn't matter in which unit the circumferences are measured, as long as it's in the same units for both of them.}}{circumference\ hip [cm/in]} = \frac{88\ cm}{103\ cm} = 0.85$
\end{fminipage}
\end{center}

\vspace{2mm}

\noindent For women, the WHR should be less than 0.85 and for men less than 1.0.
Or expressed in a more visual way: the ``pear shape'' as an external appearance is more healthy than the ``apple shape'' (see picture~\ref{Apple}).

\begin{figure}[htb!]
  \centering
  \includegraphics[width=4cm]{ApplePear}
  \caption{Apple and pear shape}\label{Apple}
\end{figure}

Newer research showed, that the waist circumference alone is the better measure to judge the abdominal fat pads (inside the belly), than the ``waist--hip--ratio''.

Gauging the health risk based on the waist circumference.

\begin{table}[htb!]
  \centering
  \begin{tabular}{lllll}
    \multicolumn{4}{c}{\textbf{Waist circumference}} & \textbf{Health risk} \\
    \textbf{Men} & & \textbf{Women} \\
    \hline
    $\leq$ 94 cm & (37 in) & $\leq$ 80 cm & (31\sfrac{1}{2}\ in) & normal \\
    95 -- 102 cm & (37 -- 40\sfrac{1}{4}\ in) & 80 --88 cm &
                  (31\sfrac{1}{2}\ -- 34\sfrac{1}{2}\ in) & elevated \\
    $>$ 102 cm & (40\sfrac{1}{4}\ in) & $>$ 88 cm  & (34\sfrac{1}{2}\ in) & strongly elevated \\
  \end{tabular}
\end{table}

To sum everything up, it can be said, that the heavier somebody is, the longer the overweight problem already consists
and the more the fat deposits are distributed around the waist, the bigger the associated health risk.

\vspace{5mm}
\noindent
\begin{fminipage}{\textwidth}
  \textbf{Profile Normal Weight}
  \begin{itemize}
  \item The Body--Mass--Index BMI should be between 18.5 -- 24.9 and with increasing age more on the upper end of this interval.
  \item The waist circumference is a god measure of the fat distribution in the body. It should be not more than 94 cm (37 in) for men and not more than 80 cm (31\textonehalf\ in) for women.
  \item Overweight most often doesn't only indicate wrong nutrition, but also a lack of movement.
    \item Determining the weight once  a month in the morning on empty stomach allows a timely correction of the behavior in terms of food and movement. 
  \end{itemize}
  \end{fminipage}

\end{document}